\scanpage{96}
\mspage{102}

c) Ces conditions restent satisfaites après tout changement de base $Y'\to Y$.

2°) Notion duale de fibration homologique complète.

\uline{Exemple} Si $f$ est $\mathrm{Hot}$-propre, les \deleted{conditions}\inserted{flèches de} a) et b) sont des $\mathrm{Hot}$-équivalences, donc en vertu de Der 3 b'$_1$ (qu'il faut admettre) ce sont des $\mathbb D$-équiv. coh\unsure{om.|om}, \deleted{\unsure{ceci|cela}} signifie que a) et b) sont valables, et c) aussi car la notion de flèche $\mathrm{Hot}$-propre est stable par chgt. de base.

\uline{Prop 2} \deleted{\unsure{Soit $t$|Soient}} $\bar f:X\unsure{\relbar|\to}Y$ $\mathbb D$-fibration coh. complète, $F\in\mathrm{Ob}\,\mathbb D(X)$ loc constant sur $X$. Alors $f_*(F)$ est loc. constant, sa formation commute à tout chgt. de base, et en particulier se calcule \deleted{\unsure{fib}}\unsure{fibre|fibres} par fibres.

% slanted
\margin{il y a une réciproque cf prop 3 !}

\uline{Cor. 1} Soit \deleted{\unsure{alors|un}} un diagramme cartésien

\diagram{\begin{tikzcd}[ampersand replacement=\&,column sep=large]
{} \& X\arrow[d,"f"] \& X'\arrow[l,"q"']\arrow[d,"f'"]\\
\deleted{\unsure{Y'|Y}}\arrow[r,dash] \& Y \& Y'\arrow[l,"\unsure{p|g}"]
\end{tikzcd}}

On suppose $f$ $\mathbb D$-fibration cohomologique, et $p$ une $\mathbb D$-équiv. coh. Alors $q$ est une $\mathbb D$-équiv. coh.

% archive sheet number: 95
